\documentclass[oneside]{ausarbeitung}
\usepackage{dirtree}
\addbibresource{latexlit.bib}
% Vermeidet einen bekannten hyperref/listings-Konflikt (Absturz "Argument of ...
% has an extra }"), wenn ein Seitenumbruch mitten in eine lstlisting-Umgebung
% fällt -- die PDF-Seitenlabels sind rein kosmetisch und werden hier nicht benötigt.
\hypersetup{pdfpagelabels=false}

\begin{document}

%--- Sprachauswahl
% Erlaubte Werte:
% \selectlanguage{english}
% \selectlanguage{ngerman}
\selectlanguage{ngerman}

%--- Art der Arbeit
% Erlaubte Werte:
% \Praxissemesterbericht
% \Projektbericht
% \Bachelorarbeit
% \Seminararbeit
% \Masterarbeit
% \Kursdokumentation{Zu dokumentierender Kurs}
\Kursdokumentation{Spieleprogrammierung}


%--- Studiengang:
% Erlaubte Werte:
% \Informatik
% \Elektronik
% \DataScience
% \DigProdDes
\DigProdDes

\title{Roll of Fate}

\author{Fabian Lux}
\matrikelnr{3005309}

%--- Ist der Erstbetreuer (\examinerA) an der Hochschule ein Professor?
% Erlaubte Werte:
% \examinerIsAProfessortrue % Ja
% \examinerIsAProfessorfalse % Nein
\examinerIsAProfessorfalse
\examinerA{Stefan Wehrenberg}
% \examinerB{Prof. Dr. Markus Weinberger}

\date{15.04.2026}

%--- Angaben zur Firma
% Auskommentieren, wenn die Arbeit nicht bei einer ext. Firma gemacht wurde.
% \companyname{Beispielfirma}
% \industrialsector{Beispielbranche}
% \department{Beispielabteilung}
% \companystreet{Beispielstr. 1}
% \companycity{12345 Musterstadt}

%--- Angaben zum Betreuer bei dieser Firma
% \advisorname{Name des Betreuers}
% \advisorphone{(01234) 567-890}
% \advisoremail{name@company.xxx}

%--- KI-Nutzungserklärung
% Bausteine per true/false an-/abschalten, beliebig kombinierbar (siehe
% ausarbeitung.cls für alle verfügbaren Schalter):
\AIsparringtrue      % KI als Sparringpartner für Konzeption & Architektur
\AItutortrue         % KI als Tutor für eine neue Technologie/Engine
\AItutortech{Godot}  % welche Technologie (nur relevant, wenn \AItutortrue)
% \AIcodetrue        % KI zur Code-Generierung
\AIdebugtrue       % KI zur Fehlersuche/Debugging (Analyse & Lösungsvorschläge, keine Code-Generierung)
% \AItexttrue        % KI zur sprachlichen Überarbeitung (Rechtschreibung/Grammatik)
% \AIresearchtrue    % KI zur Recherche/Literatursuche
% \AItranslationtrue % KI zur Übersetzung von Textabschnitten
% \aiusageentry{Eigener zusätzlicher Punkt}

%--- Signaturzeilen (Eidesstattliche Erklärung, KI-Nutzungserklärung, Sperrvermerk)
% Ohne Angabe bleibt es bei einer leeren Zeile zum handschriftlichen Ausfüllen.
% Bilder (PNG) liegen unter ./images/, werden exakt über der jeweiligen
% Beschriftung eingefügt (kein Float):
% \locationdateimage{ort_datum.png}
% \studentsignatureimage{unterschrift.png}

\maketitle
\cleardoublepage

\pagenumbering{roman}
\setcounter{page}{1}

% \makeworkplace
% \cleardoublepage

\makeaffirmation
\cleardoublepage

\makeaiclause
\cleardoublepage

% \makeconfidentialclause
% \cleardoublepage

\begin{abstract}
\enquote{Roll of Fate} ist ein rundenbasiertes Taktik-Rogue-like für mobile Plattformen, das in Godot~4.7
(GDScript) im Rahmen des Kurses Spieleprogrammierung entwickelt wurde. Kernmechanik ist ein sechsseitiger
Würfel, der sich pro Zug um 90 Grad über ein Rasterfeld rollt; die jeweils obenliegende Würfelseite
entspricht einer aktiven Fähigkeit. Über ein On-the-fly-Deckbuilding-System sammelt der Würfel im Laufe
eines Runs neue Fähigkeiten von prozedural generierten Etagen ein, während vier Gegnertypen mit
vorhersehbaren, aber tödlichen Mustern den Abstieg erschweren.

Architektonisch folgt das Spiel dem Node-/Autoload-Prinzip von Godot: GameManager und LevelManager bilden
als zentrale Zustandsmaschine bzw.\ Rasterverwaltung das Rückgrat des Spiels, während Würfelseiten und
Gegner als austauschbare, eigenständige Szenen mit fester Schnittstelle realisiert sind. Ca.\ eine Woche
vor der Abschlusspräsentation wurde das ursprünglich reine Singleplayer-Konzept um einen lokalen
WLAN-Mehrspieler-Modus erweitert, der über ein schlankes, selbstgebautes host-autoritatives Protokoll
realisiert wurde.

Die vorliegende Dokumentation beschreibt Spielkonzept, Architektur und Implementierung des Basis-Spiels
sowie die nachträgliche Koop-Erweiterung getrennt voneinander, um sichtbar zu machen, welche Anpassungen
an bereits bestehenden Systemen für die Erweiterung notwendig waren. Ergänzt wird sie um Angaben zu den
selbst erstellten 3D-Assets, zur Musiklizenz sowie zur eigens entwickelten, headless ausführbaren
Test-Suite.
\end{abstract}

\cleardoublepage
\tableofcontents

\listofabbreviations
\begin{acronym}[KI]     % Längstes Kürzel in der nachfolgenden
                        % Liste um die Breite der Spalte für die
                        % Abkürzungen zu bestimmen.
    % \acro{ki}[KI]{Künstliche Intelligenz}
    % \ac{ki} im text
\end{acronym}


\cleardoublepage
\pagenumbering{arabic}
\setcounter{page}{1}

% -----------------------------------------------------------------------
\chapter{Einleitung}
\label{cha:einleitung}

\enquote{Roll of Fate} entstand im Rahmen des Kurses \textit{Spieleprogrammierung} als rundenbasiertes
Taktik-Rogue-like für Mobilgeräte, in dem ein rollender Würfel als Spielfigur dient. Ursprünglich war
das Projekt ausschließlich als Singleplayer-Erfahrung geplant und umgesetzt; die Kapitel~\ref{cha:spielkonzept}
bis~\ref{cha:assets} beschreiben das Spiel entsprechend in diesem ursprünglichen Umfang.

Ca. eine Woche vor der Abschlusspräsentation wurde zusätzlich die Entscheidung getroffen, einen
lokalen Mehrspieler-Modus zu ergänzen. Da dieser nicht Teil der ursprünglichen Planung war, sondern
eine bewusste, nachträgliche Erweiterung darstellt, wird er getrennt von der übrigen Dokumentation in
den Kapiteln~\ref{cha:coop-konzept} und~\ref{cha:coop-architektur} behandelt.

Da Godot für den Autor zu Beginn eine unbekannte Engine war, dienten neben KI-gestützter Unterstützung
(siehe KI-Nutzungserklärung) vor allem die offizielle Dokumentation~\cite{godot:docs} sowie
Community-Tutorials, u.\,a.\ das \enquote{Godot MEGA Tutorial} von Brackeys~\cite{brackeys:megatutorial}
sowie diverse Artikel von GDQuest~\cite{gdquest:godot}, als
zentrale Lernressourcen.

\chapter{Spielkonzept}
\label{cha:spielkonzept}
\section{Überblick \& Genre}

\textbf{Genre:} Rundenbasiertes Taktik-Rogue-like mit \enquote{On-the-fly Deckbuilding}.

\textbf{Spielstruktur:} Endless Descent (Endloser Abstieg in Raster-Räumen, ähnlich wie Hoplite).

\textbf{Plattform \& Steuerung:} Mobile-fokussiert (Hochformat-Viewport, skaliert über Godots
Multiple-Resolutions-Handling~\cite{godot:multipleresolutions}); Swipe-Steuerung für diagonale Bewegung.

\textbf{Perspektive \& Look:} Obwohl die Steuerung über scheinbar diagonale Wischgesten erfolgt, operiert die Spiellogik auf einem orthogonalen Raster. Die Kamera wird um 45 Grad auf der Y-Achse rotiert, wodurch visuell eine isometrische Perspektive entsteht. Ein Wisch nach rechts-oben entspricht logisch einer Bewegung entlang der positiven Z-Achse (Forward)..

\textbf{Core-Mechanik:} Der Spieler steuert einen 6-seitigen Würfel durch ein Raster. Jede Bewegung rollt den Würfel um 90 Grad auf das Nachbarfeld. Die Fähigkeit, die nach dem Rollen oben liegt wird aktiviert. Räumliches Vorausplanen der Würfelrotation ist der Schlüssel zum Sieg.

\begin{figure}[h]
    \centering
    \fbox{\parbox[c][6cm]{10cm}{\centering TODO: Screenshot -- isometrische Ansicht des Spiels mit Würfel auf dem Raster (Kamera-Look \& Rasterbewegung)}}
    \caption{Isometrische Ansicht der Core-Mechanik}
    \label{fig:core-mechanik}
\end{figure}


\section{Der Core-Loop}
\textbf{Kampf-Räume:} Der Spieler startet auf einer prozedural generierten Ebene mit Hindernissen und Gegnern. Ziel ist es, den Abstieg zum nächsten Level zu erreichen und so möglichst viele Ebenen zu durchlaufen.

\textbf{Runden-System:} Der Kampf verläuft strikt rundenbasiert.
\begin{itemize}
    \item \textbf{Player Turn:} Spieler bewegt sich, Würfel rollt 1 Feld, Fähigkeit wird ausgelöst.
    \item \textbf{Enemy Turn:} Gegner bewegen sich oder greifen an.
\end{itemize}

\textbf{Kampf-Decks:} In manchen Räumen liegen neue \enquote{Würfelseiten} (Fähigkeiten) auf dem Boden. Rollt der Spieler darüber, wird die Würfelseite an der untenliegenden Seite angebracht und ersetzt die bisherige Würfelseite. So soll der Spieler seinen Würfel im laufe des Spiels aufbauen. Der Würfel startet mit \enquote{leeren} Würfelseiten.

\textbf{Punkte \& Heilung:} Das Erreichen einer Etage gibt 200 Punkte, ein Gegner-Kill durch tatsächlichen Schaden (z.B. per Laser-Turret, Flammenwerfer oder Dash) 30 Punkte; ein erlittener Treffer zieht 10 Punkte ab, der Punktestand fällt dabei nie unter 0. Ein Kill durch Schaden heilt zusätzlich 1 Trefferpunkt (bis zum Maximum). Das Hineinrollen in einen Schleim-Würfel dagegen tötet diesen zwar ebenfalls, zählt aber bewusst nicht als Schadens-Kill und gibt daher weder Punkte noch Heilung. Der höchste je erreichte Punktestand wird geräteweise als Highscore gespeichert.

\section{Spieler-Fähigkeiten}
Aktionen sind an spezifische Trigger-Zeitpunkte gebunden und können vor oder nach dem ausführen einer Bewegung getriggert werden oder die Bewegung selbst modifizieren.

\begin{infoblock}
    \eintrag{Flammenwerfer\newline(Offensiv)}{Verursacht Schaden auf den 2 Feldern direkt vor dem Würfel in der aktuellen Blickrichtung. Trigger: Nach dem Rollen.}
    \eintrag{Laser-Turret\newline(Offensiv)}{Feuert Laser in alle direkt umliegenden Felder (1 Feld Reichweite, 360 Grad). Trigger: Nach dem Rollen.}
    \eintrag{Holo-Schild\newline(Defensiv)}{Erzeugt eine Barriere, die den Würfel im gegnerischen Zug vor Schaden schützt. Trigger: Aktiviert nach dem Rollen.}
    \eintrag{Phase\newline(Movement)}{Der Würfel rollt nicht, sondern gleitet in Laufrichtung bis zu 3 Felder weit durch alles hindurch - Wände, Säulen und Gegner blockieren den Weg dazwischen nicht. Die Landeposition muss jedoch begehbar sein: von 3 absteigend wird die größte Distanz gewählt, deren Zielfeld kein Hindernis ist, sodass der Würfel nie in einer Wand oder Säule stehen bleibt. Dabei verschwindet der Würfel komplett und wird durch einen blauen Partikeleffekt ersetzt, der seine Position anzeigt, während er unsichtbar ist. Unabhängig von der zurückgelegten Distanz dreht er sich dabei um genau eine Seite weiter, wie bei einem einzelnen Rollschritt. Trigger: Beim Rollen.}
    \eintrag{Dash\newline(Movement)}{Der Würfel macht einen Satz und rollt sofort 2x in dieselbe Richtung (überspringt eine Seite). Gegner im Weg werden bis zum nächsten Hindernis weggeschleudert und erleiden Aufprallschaden. Trigger: Beim Rollen.}
\end{infoblock}

\section{Planungs-Modus (Ghost)}
Neben dem direkten Wisch bietet die Steuerung einen zweiten Modus, um Züge risikofrei vorab
durchzuprobieren: Hält man den Finger auf dem Würfel lange gedrückt (ausgewertet über Godots
\texttt{InputEventScreenTouch}~\cite{godot:inputeventscreentouch}), erscheint ein
halbtransparenter \enquote{Geister}-Würfel, der sich durch Ziehen des Fingers
(\texttt{InputEventScreenDrag}~\cite{godot:inputeventscreendrag}) unabhängig vom echten Würfel
Schritt für Schritt über das Raster bewegen lässt, ohne dass der reale Zug bereits ausgeführt wird.
Erst das Loslassen verwirft die Planung wieder; der tatsächliche Zug erfolgt weiterhin klassisch per
Wisch. Das erlaubt es dem Spieler zu erproben welche Würfelseite an welchem Punkt des geplanten Weges oben liegen wird.

\begin{figure}[h]
    \centering
    \fbox{\parbox[c][6cm]{10cm}{\centering TODO: Screenshot -- Long-Press mit transparentem Vorschau-Würfel neben dem echten Würfel}}
    \caption{Planungs-Modus (Ghost) während eines Long-Press}
    \label{fig:ghost-modus}
\end{figure}

\section{Anleitung für neue Spieler:innen}
Beim allerersten Spielstart blendet das Spiel Pfeile um
den Würfel ein, die auf die grundlegende Bewegung hinweisen und nach dem ersten ausgeführten Zug
automatisch verschwinden. Nach einigen weiteren Zügen folgt, sofern noch nicht selbst entdeckt, ein
zweiter Hinweis auf den oben beschriebenen Planungs-Modus. Beide Hinweise erscheinen jeweils nur
einmal pro Spielsitzung.

\begin{figure}[h]
    \centering
    \fbox{\parbox[c][6cm]{10cm}{\centering TODO: Screenshot -- Bewegungspfeile-Einblendung beim allerersten Spielstart}}
    \caption{Tutorial-Hinweis beim Einstieg}
    \label{fig:tutorial}
\end{figure}


\section{Gegner}
Gegner haben Angriffe und Bewegungsmuster, die während der Spieler zieht klar angezeigt werden. Der Spieler muss seine Endposition planen, um auszuweichen und gleichzeitig anzugreifen.

\begin{infoblock}
    \eintrag{Scharfschütze}{Bewegt sich nicht. Hat eine Reaktionszeit: Erst wenn sich der Spieler zwei Runden in Folge in dessen 4-Felder-Radius mit freier Schussbahn befindet, schießt er (Wände und Säulen blockieren). Ein einzelner Zug auf einem bedrohten Feld bleibt also folgenlos. Die Vorschau zeigt nur tatsächlich bedrohte Felder mit freier Sichtlinie, nicht den vollen Radius. Muss nach jedem Schuss 3 Runden nachladen.}
    \eintrag{Bomber}{Weicht dem Spieler aktiv aus. Lässt alle 6 Runden eine Bombe fallen. Die Bombe tickt 3 Runden lang und explodiert dann in einem 3x3 Raster, verursacht massiven Schaden an Spieler und anderen Gegnern.}
    \eintrag{Bully}{Bewegt sich jede Runde 1 Feld. Wenn seine X- oder Y-Achse mit dem Spieler übereinstimmt, stürmt er im Gegner-Zug auf den Spieler zu und schiebt ihn ans nächste Hindernis (kein direkter Schaden, der Stoß selbst ist die Gefahr - z.B. in eine Säule oder vor einen anderen Gegner). Muss nach einem Dash 3 Runden regenerieren.}
    \eintrag{Schleim-Würfel}{Bewegt sich jede 2. Runde zufällig ein Feld weiter und hinterlässt dabei eine grüne Schleimspur. Rollt der Spieler über den Schleim, wird die Seite, die den Boden berührt (die Unterseite), verklebt und bewegt sich fortan mit dem Würfel mit; sie wird erst wieder freigegeben, sobald sie erneut Bodenkontakt hat. Der Schleim-Würfel selbst verursacht keinen Schaden, ist aber gefährlich: Landet der Spieler am Ende seines Zugs auf dem Feld, auf das der Schleim-Würfel als Nächstes ziehen wird (rot markiert), oder rollt er direkt in ihn hinein, verkleben alle sechs Würfelseiten auf einmal und der Schleim-Würfel stirbt. Das Hineinrollen in jeden anderen Gegnertyp blockiert den Zug stattdessen einfach wie eine Wand.}
\end{infoblock}

\begin{figure}[h]
    \centering
    \fbox{\parbox[c][6cm]{10cm}{\centering TODO: Screenshot -- Gefahrenzonen-Vorschau eines Gegners (rote Partikel-Umrisse)}}
    \caption{Gefahrenzonen-Vorschau}
    \label{fig:gegner-vorschau}
\end{figure}

\chapter{Architektur}
\label{cha:architektur}
 
\section{Grundprinzip}
 
Die Architektur folgt dem Node-basierten Kompositionsprinzip von Godot~4~\cite{godot:engine}. Das Spiel wird in eigenverantwortliche Szenen aufgeteilt, die über Signals~\cite{godot:signals} miteinander kommunizieren. Direkte Referenzen zwischen Systemen werden vermieden, um Kopplung zu minimieren. Systemübergreifende Dienste werden als Autoloads (Godots Singleton-Mechanismus~\cite{godot:autoload}) registriert und sind so aus jeder Szene global erreichbar.
 
\section{Systeme}
 
Das Spiel gliedert sich in fünf Kernsysteme.
 
\subsection{Würfel}
 
Der Würfel besteht aus zwei getrennten Verantwortlichkeiten: dem Würfel selbst und seinen Seiten.
 
Der \textbf{Würfel-Node} verwaltet die Bewegung auf dem Raster, die visuelle Rollanimation und die Berechnung, welche Seite oben liegt. Er fragt vor jeder Bewegung beim LevelManager an, ob das Zielfeld begehbar ist, und führt die Aktionen der Würfelseiten aus.
 
Die \textbf{Würfelseiten} sind eigenständige Szenen, die als Kind-Nodes in die sechs Slot-Marker des Würfels eingehängt werden. Jede Seite besteht aus einem 3D-Modell, einem AnimationPlayer und einem Logik-Skript mit einer fest definierten Schnittstelle:
 
\begin{infoblock}
    \eintrag{\texttt{pre\_activate()}}{Wird aufgerufen wenn diese Seite zur Oberseite wird, bevor die Bewegung ausgeführt wird.}
    \eintrag{\texttt{deactivate()}}{Wird aufgerufen, wenn diese Seite die Oberseite verlässt, bevor die Bewegung ausgeführt wird.}
    \eintrag{\texttt{post\_deactivate()}}{Wird aufgerufen, wenn diese Seite die Oberseite verlassen hat und die Bewegung vorbei ist.}
    \eintrag{\texttt{activate()}}{Wird aufgerufen wenn diese Seite zur Oberseite geworden ist und die Bewegung vorbei ist.}
\end{infoblock}

Zusätzlich erlaubt \texttt{modify\_roll(player, direction) -> bool} der aktuell aktiven (obenliegenden) Seite, die normale Roll-Bewegung komplett zu übersteuern (genutzt von \textit{Phase} und \textit{Dash}). Gibt die Methode \texttt{true} zurück, entfällt der reguläre Bewegungsablauf inklusive der obigen vier Methoden für diesen Zug.
 
Manche Seiten besitzen zusätzlich ein Partikelsystem für visuelle Effekte. Da jede Seite vollständige Visualisierung und Logik in einer Szene vereint, können neue Fähigkeiten als eigenständige \texttt{.tscn}-Dateien entwickelt und ohne Änderungen am Würfel-System eingesetzt werden.
 
\subsection{LevelManager}
 
Der \textbf{LevelManager} verwaltet das Spielfeld als zweidimensionales Raster (statische Geometrie: Boden, Wand, Säule, Start, Ausgang, Pickup) und ist der einzige Ort im Spiel, an dem der aktuelle Weltzustand zu Wänden/Hindernissen abgefragt werden kann. Die zentrale Liste aller dynamischen Akteure (Spieler, Gegner, abgelegte Entities) wird dagegen vom \textbf{GameManager} geführt, da Spieler- und Gegner-Registrierung ohnehin über diesen laufen und so eine zweite Quelle für denselben Zustand vermieden wird: statischer Geometrie-Lookup im LevelManager, dynamische Entity-Liste im GameManager.

Der LevelManager übernimmt außerdem die prozedurale \textbf{Level-Generierung} zu Beginn jeder Etage. Ein untergeordneter \texttt{LevelGenerator}-Node erzeugt das Raster, platziert Hindernisse, Gegner, Würfelseiten-Pickups und den Level-Ausgang. Da die Generierung direkt auf den internen Rasterdaten arbeitet, ist eine Integration in den LevelManager sinnvoller als ein eigenständiges Top-Level-System.
 
\subsection{GameManager}
 
Der \textbf{GameManager} ist das zentrale Steuerungselement des Spiels und als Autoload implementiert. Er verwaltet den gesamten Core-Loop und ist für alle übergreifenden Zustandsübergänge zuständig. Intern arbeitet er als State Machine~\cite{gdquest:fsm}:
 
\begin{infoblock}
    \eintrag{\texttt{PLAYER\_TURN}}{Spielereingabe wird akzeptiert.}
    \eintrag{\texttt{RESOLVING}}{Fähigkeit der Oberseite und Effekte werden abgearbeitet.}
    \eintrag{\texttt{ENEMY\_TURN}}{Alle Gegner führen nacheinander ihre Aktionen aus.}
    \eintrag{\texttt{GAME\_OVER} / \newline\texttt{LEVEL\_COMPLETE}}{Endzustände einer Etage und Übergang zur nächsten \newline Etage.}
\end{infoblock}
 
Der GameManager führt nach jedem Zustandsübergang die relevanten Checks aus: Ist der Spieler gestorben? Wurde der Ausgang betreten? Würfel, Gegner und Level-Elemente kommunizieren dabei größtenteils über direkte Methodenaufrufe auf den Autoloads (\texttt{GameManager}, \texttt{LevelManager}) statt ausschließlich über Signals. Signals werden für Zustands\"ubergänge genutzt, auf die mehrere unabhängige Systeme reagieren (z.B. \texttt{player\_turn\_started}), während 1:1-Anfragen (z.B. „ist Feld X begehbar?“, „füge Spieler Schaden zu“) aus Einfachheitsgründen als direkte Aufrufe umgesetzt sind.

\subsection{Gegner}
 
Analog zu den Würfelseiten sind alle Gegnertypen als eigenständige Szenen umgesetzt. Jede Szene besteht aus einem 3D-Modell, einem AnimationPlayer und einem Logik-Skript, das von einer gemeinsamen Basisklasse erbt. Die Basisklasse definiert eine feste Schnittstelle:
 
\begin{infoblock}[0.3]
    \eintrag{\texttt{take\_turn()}}{Wird vom GameManager aufgerufen; der Gegner führt Bewegung und Aktion aus.}
    \eintrag{\texttt{take\_damage(amount)}}{Verarbeitet eingehenden Schaden.}
    \eintrag{\texttt{predict()}}{Zeigt dem Spieler an, was der Gegner im nächsten Zug tun wird.}
\end{infoblock}
 
Jeder Gegnertyp überschreibt diese Methoden mit seinem spezifischen Verhalten.

\subsection{SoundManager}

Der \textbf{SoundManager} ist als Autoload für sämtliche Audioausgabe zuständig und trennt dabei
Hintergrundmusik von Soundeffekten. Er wählt beim Start eine gemischte Musik-Playlist aus und spielt
für punktuelle Ereignisse (Bewegung, Treffer, Kills, UI-Interaktion, ...) passende Soundeffekte ab (wobei hier nur ein Lied eingefügt wurde, also keine echte Playlist).

\subsection{UI \& Hauptmenü}

Die Benutzeroberfläche umfasst zwei Bereiche. Das \textbf{Hauptmenü} bietet Spielstart, Einstellungen (Lautstärke, Grafikoptionen), einen Credits-Bildschirm sowie den bisher erreichten Highscore. Beim ersten Spielstart ergänzt ein einmaliger Tutorial-Hinweis (Bewegungspfeile, später ein Hinweis auf den Planungs-Modus) die Einführung neuer Spieler:innen. Das \textbf{In-Game-HUD} zeigt die aktuelle Etage und den Punktestand, warnt bei kritischen Trefferpunkten über eine rote Vignette und blendet einen Hinweis ein, solange das eigene Gerät nicht am Zug ist. Beide Bereiche sind als eigenständige Szenen implementiert und werden vom GameManager ein- und ausgeblendet.
 
\section{Kommunikationsfluss}
 
Der typische Ablauf eines Spielerzugs verdeutlicht das Zusammenspiel der Systeme:
 
\begin{enumerate}
    \item Der Spieler führt eine Swipe-Geste aus; der Würfel-Node empfängt die Richtung.
    \item Der Würfel-Node fragt den LevelManager, ob das Zielfeld begehbar ist.
    \item Bei gültigem Zug wird die Rollanimation abgespielt.
    \item Die bisherige Oberseite wird deaktiviert, die neue Oberseite aktiviert.
    \item Der Würfel-Node ruft \texttt{GameManager.on\_player\_turn\_ended()} auf.
    \item Der GameManager wechselt in \texttt{RESOLVING}, führt alle Checks aus und leitet danach den \texttt{ENEMY\_TURN} ein.
    \item Jeder Gegner wird der Reihe nach über \texttt{await enemy.take\_turn()} aufgerufen (bewusst direkt awaited statt über ein separates Signal, um eine Race Condition zu vermeiden: bei synchron abgeschlossenen Zügen könnte ein vorher emittiertes Signal sonst verpasst werden).
    \item Der GameManager wechselt zurück in \texttt{PLAYER\_TURN}.
\end{enumerate}

\chapter{Implementierung}
\label{cha:implementierung}

\section{Würfel}

Die größte technische Herausforderung der Würfel-Schnittstelle ist die saubere Trennung von visueller Rotation und logischer Positionsbestimmung auf dem Raster. Ein Würfel rotiert bei Fortbewegung nicht um seinen Mittelpunkt, sondern kippt über seine untere Kante. 

Um dies in der Engine umzusetzen, wird ein verschiebbarer Drehpunkt (Pivot) genutzt. Die Node-Hierarchie in Godot ist wie folgt aufgebaut:\\

\dirtree{%
.1 Player (Node3D).
.2 Pivot (Node3D).
.3 CubeMesh (MeshInstance3D).
.4 Slot\_1 bis Slot\_6 (Marker3D).
.5 [Ability-Szenen].
.2 CameraBase (Node3D).
.3 Camera3D.
}

\textbf{Ablauf einer Rotations-Phase:}
Wird ein valider Input registriert, durchläuft das Skript deterministische Phasen, um asynchrone Fehler bei Animationen zu vermeiden:

\begin{enumerate}
    \item \textbf{Vorbereitung:} Der \texttt{Pivot}-Node wird programmatisch an die Kante verschoben, über die abgerollt werden soll (halbe Grid-Größe in Bewegungsrichtung). Das \texttt{CubeMesh} wird simultan in die entgegengesetzte Richtung verschoben, um seine absolute Weltposition zunächst beizubehalten.
    \item \textbf{Animation:} Ein Godot-\texttt{Tween}~\cite{godot:tween} rotiert den Pivot-Node parallel zur Kamera-Bewegung um 90 Grad.
    \item \textbf{Cleanup \& Trigger:} Nach Abschluss der Animation wird die neue Weltposition auf den Haupt-Node (\texttt{Player}) übertragen, der Pivot wird in den Ursprung zurückgesetzt und die Rotation in die \texttt{MeshInstance} \enquote{gebacken}.
\end{enumerate}

Durch diesen Ablauf behält der Haupt-Node seine reine Raster-Position bei, während das komplexe Rotationsverhalten vollständig in den Child-Nodes gekapselt wird.

Außerdem werden in diesen Ablauf mehrere Methodenaufrufe gestartet. Diese informieren sowohl die aktuelle Oberseite, als auch die zukünftige Oberseite über die Rotation. Dadurch können zu den richtigen Zeitpunkten (vor dem Rollen und nach dem Rollen) Fähigkeiten aktiviert und deaktiviert werden.

Die Ermittlung der zukünftigen Oberseite erfolgt vor der Bewegung deterministisch über das Skalarprodukt der lokalen Marker-Ausrichtungen und dem invertierten Bewegungsvektor.

\textbf{Schnittstellen und Kapselung}
Um eine lose Kopplung zwischen den Systemen zu gewährleisten, agiert der Würfel weitestgehend selbstständig und kommuniziert über klar definierte Ein- und Ausgänge.

\begin{itemize}
    \item \textbf{Zustandsverwaltung \& Input:} Der Würfel wertet Swipe-Gesten eigenständig aus. Ein interner Status (\texttt{awaiting\_input}) stellt sicher, dass weitere Eingaben blockiert werden, solange eine Rotation oder eine Fähigkeit ausgeführt wird.
    
    \item \textbf{Deckbuilding-Schnittstelle:} Für das Modifizieren des Würfels zur Laufzeit wird die Methode \texttt{attach\_to\_bottom(new\_scene: PackedScene) -> PackedScene} bereitgestellt. Diese tauscht die Szene im aktuell untenliegenden Slot aus und gibt die verdrängte Szene zurück (z.B. eine Basis-Seite mit Würfelaugen oder eine ersetzte Fähigkeit). Dies erlaubt es dem GameManager, die abgelegte Seite als neues Objekt auf dem Raster zu instanziieren.
    
    \item \textbf{Signal-Kommunikation (Outbound):} Da die Fähigkeiten als eigenständige Szenen gehandhabt werden, reicht ein simples Ende der Roll-Animation nicht aus, um den Zug zu beenden. Die Signal-Kette ist delegiert: Nach dem Rollen ruft der Würfel \texttt{activate()} auf der neuen Oberseite auf. Sobald die Ability-Szene ihre Effekte abgeschlossen hat, meldet sie dies an den Würfel zurück, welcher daraufhin das Signal \texttt{turn\_ended} an den GameManager emittiert. Flankierend feuert der Würfel Signale wie \texttt{roll\_started(target\_pos)}, um dem System Positionswechsel frühzeitig mitzuteilen.
\end{itemize}

\section{LevelManager}

Der \texttt{LevelManager} fungiert als globale Schnittstelle für den Level-Zustand und trennt dabei strikt zwischen statischer Geometrie und dynamischen Akteuren.

\textbf{Statische Daten (Level-Layout):}
Das beim Etagen-Start generierte Raster wird in einem Dictionary gespeichert, welches \texttt{Vector2i}-Koordinaten auf \texttt{GridCell}-Objekte abbildet. Diese Zellen speichern ausschließlich passive, statische Eigenschaften (z.B. ob es sich um eine unpassierbare Wand handelt).

\textbf{Dynamische Daten (Entities):}
Um redundante Datenhaltung und die damit verbundene fehleranfällige Synchronisation zu vermeiden, speichern die \texttt{GridCells} \textit{nicht}, ob sie aktuell von einem Charakter blockiert werden. 
Stattdessen führt der \texttt{GameManager} die zentrale Liste aller generierten Gegner und des Spielers (\texttt{Entities}); jede Entity ist selbst die Single Source of Truth für ihre aktuelle Rasterposition.

Muss ein System prüfen, ob ein Feld blockiert ist, prüft es über den \texttt{LevelManager} das statische Raster auf Wände und über den \texttt{GameManager} die Entity-Liste. Bei der angestrebten geringen Anzahl von maximal ca. 20 Entities pro Etage ist dieser lineare Suchaufwand vernachlässigbar und garantiert eine absolute Datenkonsistenz.

\textbf{Auswertung durch Fähigkeiten:}
Dieses Design verschiebt die Regelauswertung vollständig in die Spielfähigkeiten. Der \texttt{LevelManager} agiert nur als Informationsquelle, nicht als Regelwerk.

\subsection{Prozedurale Level-Generierung}

Die Generierung neuer Etagen erfolgt zu Beginn jedes Levels deterministisch in einer mehrstufigen Pipeline. Der Algorithmus setzt auf eine additive Rechteck-Struktur, die organische, aber Raster-konforme Level-Layouts erzeugt.

\textbf{1. Layout:}
Pro Etage wird additiv ein großer, länglicher Hauptraum (kurze Seite 4-8, lange Seite 8-14 Kacheln, zufällig horizontal oder vertikal orientiert) sowie 4-6 weitere, kleinere Räume (Seitenlänge 4-7 Kacheln) erzeugt. Jeder weitere Raum wird an einen zufälligen, bereits platzierten Raum angehängt und über einen 1 breiten Korridor mit ihm verbunden. Mit einer gewissen Wahrscheinlichkeit erhalten zwei geometrisch benachbarte, aber nicht direkt verbundene Räume einen zusätzlichen Korridor, sodass gelegentlich Etagen mit mehr als einem Laufweg entstehen. Außerdem kann eine kleine Sackgasse (2x2 bis 3x3) an einen zufälligen Raum angehängt werden; ihr Mittelpunkt ist reserviert für ein Würfelseiten-Spawn (siehe Schritt 4). Diese schmalen Passagen zwingen den Spieler, sich mit Gegnern auf engstem Raum zu konfrontieren oder Fortbewegungsfähigkeiten wie \textit{Phase} taktisch einzusetzen, um blockierte Felder zu überspringen.

\textbf{2. Wände:}
Sobald die finale Bodenstruktur berechnet ist, generiert das System die unpassierbaren Außenwände. Der Algorithmus iteriert über alle validen Boden-Kacheln und prüft deren acht direkte und diagonale Nachbarn. Jedes Nachbarfeld, das noch keine Zuweisung besitzt, wird in eine Wand-Kachel umgewandelt. Durch die in Schritt 1 definierten Korridore entstehen in dieser Phase automatisch massive Wandstrukturen, die die Räume klar voneinander abgrenzen.

\textbf{3. Säulen:}
Um allzu große, offene Kampfflächen aufzubrechen und Deckungsmöglichkeiten zu schaffen, wird ein entfernungsbasierter Algorithmus eingesetzt. Für jede Boden-Kachel wird geprüft, ob sich im Umkreis einer Manhattan-Distanz von bis zu 3 Kacheln bereits eine Wand oder eine etablierte Säule befindet. Ist dies der Fall, bleibt das Feld frei. Wird jedoch in diesem Radius kein Hindernis gefunden, wird das Feld in eine Säule umgewandelt. Diese Regel sorgt für eine gleichmäßige, organische Verteilung von Deckungselementen in ausufernden Raumstrukturen.

\textbf{4. Spawn-System:}
Nachdem das Layout generiert ist, werden die dynamischen Objekte platziert:

\begin{infoblock}
    \eintrag{Start \& Ausgang}{Der Spawn-Punkt des Spielers und der Aufstieg zur nächsten Etage werden an den zwei (annähernd) am weitesten voneinander entfernten Boden-Kacheln platziert, um einen langen Weg durch das Level zu erzwingen.}
    \eintrag{Gegner}{Die Etage verfügt über ein Punktebudget von $4 + 2 \times \text{Etage}$ Punkten. Gegnertypen besitzen feste Kosten: Schleim-Würfel 3, Scharfschütze 3, Bully 4, Bomber 5 Punkte. Der Algorithmus wählt zufällig unter den noch bezahlbaren, freigeschalteten Gegnertypen, bis kein weiterer Gegner mehr ins Budget passt, und platziert diese auf zufälligen freien Boden-Kacheln außerhalb des Start-Raums.

    In den ersten vier Ebenen werden die Gegnertypen in der Reihenfolge Schleim-Würfel, Scharfschütze, Bully, Bomber nacheinander freigeschaltet (Etage $n$ schaltet die ersten $n$ Typen frei), ab Ebene 4 sind alle verfügbar.}
    \eintrag{Würfelseiten}{Existiert genau dann, wenn in Schritt 1 eine Sackgasse erzeugt wurde, und liegt an deren Mittelpunkt - die Sackgasse ist also die ausschließliche Spawn-Position für Würfelseiten. Die konkrete Würfelseite wird zufällig aus dem Pool der implementierten Fähigkeiten gezogen (Laser-Turret, Holo-Schild, Phase, Dash, Flammenwerfer).}
\end{infoblock}

\section{Sound}

Der \texttt{SoundManager} trennt Hintergrundmusik und Soundeffekte technisch voneinander.

\textbf{Soundeffekte:} Um zu verhindern, dass schnell aufeinanderfolgende Ereignisse (z.B. mehrere
Treffer kurz hintereinander) sich gegenseitig abschneiden, unterhält der SoundManager einen Pool
mehrerer \texttt{AudioStreamPlayer}-Nodes und verteilt eingehende Soundeffekte im Round-Robin-Verfahren
auf diesen Pool, statt einen einzelnen Player wiederzuverwenden.

\textbf{Musik:} Die Hintergrundmusik wird als gemischte Playlist abgespielt. Die einzelnen Titel
folgen dem Namensschema \texttt{game\_music\_NN.mp3} und werden nicht per Verzeichnis-Scan
eingesammelt, sondern über eine feste, fortlaufend hochgezählte Nummer mittels
\texttt{ResourceLoader.exists()} geprüft und geladen. Diese Umsetzung wurde bewusst gewählt, da ein
Verzeichnis-Scan sich in gepackten Android-Exports~\cite{godot:androidexport} als unzuverlässig erwiesen hat.

% \section{Tests}

% Für das Projekt existiert kein Godot-Test-Addon; stattdessen wurde ein minimales, eigenes
% Test-Framework umgesetzt (\texttt{check()}/\texttt{check\_eq()} als einfache Pass/Fail-Sammlung). Die
% Tests laufen headless über
% \begin{lstlisting}
% godot --headless --path . scenes/tests/test_runner.tscn
% \end{lstlisting}
% und liefern einen CI-tauglichen Exit-Code (0 = alle Tests bestanden, 1 = mindestens ein Fehlschlag).

% Anders als klassische Unit-Tests laufen die Tests integrationsartig gegen eine echte, prozedural
% generierte Level-Instanz und reale Node-Objekte statt gegen Mocks; da Level-Layouts jedes Mal neu
% generiert werden, suchen manche Tests aktiv nach einem passenden lokalen Layout (z.B. einer freien
% geraden Linie) oder erzwingen es, statt feste Koordinaten anzunehmen.

% Grob lassen sich die rund 20 Testfälle folgenden Kategorien zuordnen: Level-Generierung
% (Erreichbarkeit, Spawn-Budget), Bewegungs-/Kollisionsregeln, Deckbuilding, jede der fünf
% Würfelseiten-Fähigkeiten einzeln, jeder der vier Gegnertypen einzeln, sowie gezielte
% Regressionstests für zuvor gefundene Bugs (z.B. übersprungene Gegnerzüge bei einem Tod mitten in der
% Iteration, oder nicht bereinigte Gefahrenzonen-Markierungen bei einem Etagenwechsel).

% \begin{figure}[h]
%     \centering
%     \fbox{\parbox[c][6cm]{10cm}{\centering TODO: Screenshot/Konsolenausgabe eines vollständigen, erfolgreichen Testlaufs}}
%     \caption{Erfolgreicher headless Testlauf}
%     \label{fig:testlauf}
% \end{figure}


\chapter{Assets \& Content}
\label{cha:assets}

\section{3D-Assets}

Sämtliche 3D-Modelle (Würfel, Würfelseiten, Gegner, Effekt-Objekte) wurden selbst in Blender erstellt
% (Quelldateien z.B. \texttt{die.blend}, \texttt{slime.blend}, \texttt{Face\_Flamethrower.blend},
% \texttt{Face\_Shield.blend}) 
und als \texttt{.glb} exportiert .
% (u.a. \texttt{die.glb},
% \texttt{die\_face\_1.glb} bis \texttt{die\_face\_6.glb}, \texttt{face\_flamethrower.glb},
% \texttt{face\_shield\_generator.glb}, \texttt{slime.glb}, \texttt{turret.glb})


Als einheitliches Gestaltungsprinzip folgen die Modelle einem \enquote{Voxel-y}-Look: Eine Rastereinheit
entspricht einem Block von $7 \times 7 \times 7$ Voxeln. Einzelne Modelle weichen davon jedoch gezielt ab, wo ein starres Voxel-Raster ein Element nicht sinnvoll abbilden könnte. Ein Beispiel sind die Kanonen des Turret-Moduls, die kleinere Quader benötigen, als sich mit dem Voxel-Raster direkt umsetzen ließen.

\begin{figure}[h]
    \centering
    \fbox{\parbox[c][6cm]{10cm}{\centering TODO: Blender-Screenshot eines Modells (z.B. Würfel oder Slime) mit sichtbarem 7x7x7-Voxel-Raster}}
    \caption{Voxel-Raster als Gestaltungsgrundlage}
    \label{fig:voxel-raster}
\end{figure}

\begin{figure}[h]
    \centering
    \fbox{\parbox[c][6cm]{10cm}{\centering TODO: Vergleichsbild eines strikt voxeligen Modells neben dem Turret-Modul (Kanonen weichen vom Raster ab)}}
    \caption{Bewusste Abweichung vom Voxel-Raster am Beispiel des Turret-Moduls}
    \label{fig:voxel-abweichung}
\end{figure}

\section{Musik}

Die Hintergrundmusik und Soundeffekte stehen unter der Lizenz CC BY 4.0. Autor ist Dominik Braun (\url{https://dominik-braun.net}). Autor und Lizenz werden im Credits-Panel des Hauptmenüs verlinkt genannt.

\section{Partikeleffekte}

Die visuellen Effekte des Spiels (Gefahrenzonen-Vorschau, Ability-Effekte) basieren technisch auf
Godots \texttt{GPUParticles3D}~\cite{godot:particles} in Kombination mit \texttt{ParticleProcess}-\texttt{Material}~\cite{godot:particleprocessmaterial}.

Ein anschauliches Beispiel ist die Gefahrenzonen-Vorschau des Snipers: Pro Kante der markierten
Zone kommen zwei Partikelsysteme gleichzeitig zum Einsatz: Eines zeichnet dauerhaft den Umriss der
Kante als Linie, das andere lässt fortlaufend Partikel in Richtung der Quelle der Gefahr zurückfliegen, um Ursache und Wirkung visuell miteinander zu
verknüpfen. Da \texttt{ParticleProcessMaterial.direction} trotz globaler Koordinaten relativ zur
aktuellen Rotation des jeweiligen Emitters interpretiert wird, muss die gewünschte Weltrichtung vorher
in den (bei manchen Kanten um 90 Grad gedrehten) lokalen Emitter-Raum zurückgerechnet werden.

Beim Übergang zwischen zwei normalen Spielzügen laufen bereits gespawnte Partikel sanft aus, während eine Level-Regeneration alle Vorschau-Partikel sofort entfernt, da ein langsames Ausblenden über einen Etagenwechsel hinweg optisch irritieren würde.

\begin{figure}[h]
    \centering
    \fbox{\parbox[c][6cm]{10cm}{\centering TODO: In-Game-Nahaufnahme einer Gefahrenzonen-Kante mit den zwei Partikelsystemen}}
    \caption{Partikel-Nahaufnahme einer Sniper-Gefahrenzone}
    \label{fig:partikel-nahaufnahme}
\end{figure}

Weitere Ability-Effekte (Flammenwerfer, Turret, Schild, gerichtete Partikelströme) nutzen dieselbe
technische Basis mit jeweils eigens abgestimmten \texttt{Particle}-\texttt{ProcessMaterial}-Einstellungen.


\chapter{Erweiterung: Lokaler Mehrspieler}
\label{cha:coop-konzept}

\section{Motivation \& Zeitpunkt}

\enquote{Roll of Fate} wurde ursprünglich ausschließlich als Singleplayer-Spiel konzipiert und über weite
Teile des Kurses auch so umgesetzt. Alle in den Kapiteln~\ref{cha:spielkonzept}
bis~\ref{cha:assets} beschriebenen Systeme entstanden ohne Rücksicht auf einen zweiten Spieler. Ungefähr eine Woche vor der Abschlusspräsentation fiel die bewusste Entscheidung, das Spiel zusätzlich um
einen lokalen Mehrspieler-Modus zu erweitern. Dieser war nicht Teil der ursprünglichen Planung,
sondern eine nachträgliche, in begrenzter Zeit umgesetzte Ergänzung.

\section{Spielerlebnis zu zweit}

Zwei Geräte im selben WLAN können sich zu einem gemeinsamen Run verbinden, ganz ohne Internet oder
externen Server. Ein Gerät übernimmt die Rolle des \textbf{Host}, das andere die des
\textbf{Secondary}. Der Host generiert das Level; der Secondary tritt über eine Liste gefundener
Hosts bei und erhält anschließend exakt dasselbe Level-Layout wie der Host, statt es lokal separat
neu zu würfeln.

\begin{figure}[h]
    \centering
    \fbox{\parbox[c][6cm]{10cm}{\centering TODO: Screenshot des Koop-Panels im Hauptmenü (Hosten/Beitreten, Geräteliste, Namenswahl)}}
    \caption{Koop-Panel im Hauptmenü}
    \label{fig:koop-panel}
\end{figure}

\section{Rundenablauf zu zweit}

Der aus Kapitel~\ref{cha:spielkonzept} bekannte Runden-Ablauf wird um einen zweiten Spielerzug
erweitert: Auf den Zug des Hosts folgt der Zug des Secondary, erst danach sind die Gegner an der
Reihe. Der Kernloop lautet also \texttt{PLAYER\_TURN} (Host) $\rightarrow$
\texttt{SECONDARY\_TURN} $\rightarrow$ \texttt{ENEMY\_TURN} $\rightarrow$ \texttt{PLAYER\_TURN}.

\section{Sieg-/Ausgangsbedingung}

Da beide Spieler:innen gemeinsam absteigen, reicht es im Koop-Modus nicht, dass eine Person den
Levelausgang erreicht -- beide Würfel müssen sich auf einem Ausgangsfeld befinden, damit die Etage
als abgeschlossen gilt und der Übergang zur nächsten Ebene erfolgt.

\section{Interaktion zwischen den Würfeln}

Die beiden Würfel können sich gegenseitig nicht durchdringen: Der jeweils andere Spieler-Würfel
blockiert ein Feld wie ein Hindernis. Das ist insbesondere für den Dash relevant, der ohne diese
Regel durch den anderen Würfel hindurchrutschen würde, statt an ihm gestoppt zu werden.


\chapter{Architektur \& Umsetzung des Mehrspieler-Modus}
\label{cha:coop-architektur}

\section{Das Netzwerksystem (NetworkManager)}

Als neues Autoload kommt der \textbf{NetworkManager} hinzu. Bewusst wird dabei \textit{kein}
Godot-Multiplayer-Framework~\cite{godot:multiplayerdocs} eingesetzt (\texttt{ENetMultiplayerPeer}/RPCs).
Als Lernressource für den generellen Ablauf diente dabei zunächst u.\,a.\ das Tutorial \enquote{Add
Multiplayer to your Godot Game!} von Brackeys~\cite{brackeys:multiplayer}; dessen ENet-basierter Ansatz
wurde für \enquote{Roll of Fate} jedoch bewusst nicht übernommen, da er für eine feste 1:1-Verbindung
mit host-autoritativer Simulation mehr Overhead mitbringt als nötig. Stattdessen implementiert der
NetworkManager ein schlankes, selbstgebautes Protokoll direkt auf Basis von Godots
\texttt{TCPServer}~\cite{godot:tcpserver} und \texttt{StreamPeerTCP}~\cite{godot:streampeertcp}:

\begin{infoblock}
    \eintrag{Discovery}{Der Host sendet im Sekundentakt einen UDP-Broadcast~\cite{godot:packetpeerudp}
    mit seinem Namen und
    TCP-Port. Ein Host, von dem 3 Sekunden lang kein neues Paket ankam, gilt als weg -- es gibt kein
    explizites Paket für \enquote{Host beendet}, daher rein zeitbasiertes Timeout.}
    \eintrag{Verbindungsaufbau}{Nach Auswahl eines Hosts baut der Secondary eine TCP-Verbindung auf.
    Beide Seiten tauschen unmittelbar danach ihren Anzeigenamen über eine Nachricht vom Typ \enquote{hello} aus.}
    \eintrag{Nachrichtenformat}{Da TCP ein reiner Byte-Strom ohne Nachrichtengrenzen ist, bekommt
    jede Nachricht einen 4-Byte-Längenpräfix vorangestellt, über den sie beim Empfänger wieder exakt
    aus dem Datenstrom herausgeschnitten werden kann.}
\end{infoblock}

\section{Protokoll-Schemas}

Alle Nachrichten sind flache Dictionaries mit einem \texttt{type}-Feld. Die folgenden Schemas zeigen
die tatsächlich verwendeten Felder je Nachrichtentyp.

\shorthandoff{"}

\textbf{host\_announce} (UDP, Host $\rightarrow$ Suchende, periodisch jede Sekunde):
\begin{lstlisting}
{
  "type": "host_announce",
  "name": string,       // Anzeigename des Hosts
  "tcp_port": int
}
\end{lstlisting}

\textbf{hello} (TCP, erste Nachricht nach Verbindungsaufbau, in beide Richtungen):
\begin{lstlisting}
{
  "type": "hello",
  "name": string        // Anzeigename des Absenders
}
\end{lstlisting}

\textbf{turn\_request} (TCP, Secondary $\rightarrow$ Host, bei eigenem Zugwunsch):
\begin{lstlisting}
{
  "type": "turn_request",
  "direction": Vector3  // gewünschte Roll-Richtung
}
\end{lstlisting}

\textbf{turn\_action} (TCP, Host $\rightarrow$ Secondary, nach jedem angenommenen Zug):
\begin{lstlisting}
{
  "type": "turn_action",
  "actor": "host" | "secondary",
  "direction": Vector3,
  "hp": int,
  "score": int,
  "player_shielded": bool
}
\end{lstlisting}

\textbf{turn\_rejected} (TCP, Host $\rightarrow$ Secondary, kein weiterer Payload):
\begin{lstlisting}
{
  "type": "turn_rejected"
}
\end{lstlisting}

\begin{minipage}{\linewidth}
\textbf{enemy\_turn\_result} (TCP, Host $\rightarrow$ Secondary, nach jeder Gegnerrunde -- ein
vollständiger Schnappschuss statt eines reinen Bewegungs-Diffs, da ein einzelner Gegnerzug beliebig
viele Nebenwirkungen haben kann, z.B. eine neu abgeworfene Bombe oder eine mitgetötete Nachbar-Einheit):
\begin{lstlisting}
{
  "type": "enemy_turn_result",
  "hp": int,
  "score": int,
  "player_shielded": bool,
  "host_pos": Vector3i,
  "secondary_pos": Vector3i,
  "enemies": [
    {
      "sync_id": int,          // stabile ID statt Array-Index
      "sync_type": string,     // z.B. "slime", "sniper", "bully", "bomber"
      "pos": Vector3i,
      "hp": int,
      "cooldown_current": int,
      // je nach sync_type zusätzlich, z.B.:
      "next_move": Vector3i,   // Slime
      "trace_at": Vector3i,    // Slime
      "ticks_remaining": int,  // Bomb
      "fired": bool            // Sniper
    }
  ]
}
\end{lstlisting}
\end{minipage}

\textbf{level\_data} (TCP, Host $\rightarrow$ Secondary, einmalig beim Levelstart):
\begin{lstlisting}
{
  "type": "level_data",
  "floor": int,
  "cells": { Vector2i: int },      // Rasterposition -> Zelltyp
  "start_position": Vector2i,
  "secondary_start_position": Vector2i,
  "exit_position": Vector2i,
  "exit_position_2": Vector2i,
  "pickups": [ ... ],
  "enemies": [ ... ]
}
\end{lstlisting}

\textbf{state\_sync} (TCP, Host $\rightarrow$ Secondary, bei Zustandswechseln wie \texttt{GAME\_OVER}
oder \texttt{LEVEL\_COMPLETE}):
\begin{lstlisting}
{
  "type": "state_sync",
  "state": int   // Wert aus dem State-Enum des GameManagers
}
\end{lstlisting}

\shorthandon{"}

\section{Notwendige Anpassungen an bestehenden Systemen}

Der Koop-Modus wurde nicht als eigenständiges Extra-System neben das bestehende Spiel gestellt,
sondern erforderte gezielte Anpassungen an mehreren bereits bestehenden Systemen aus
Kapitel~\ref{cha:architektur} und~\ref{cha:implementierung}:

\begin{infoblock}
    \eintrag{GameManager}{Die State Machine wurde um den zusätzlichen Zustand
    \texttt{SECONDARY\_TURN} erweitert; sämtliche zugrelevante Entscheidungslogik prüft zusätzlich
    die aktuelle \texttt{NetworkManager.role}. Für den zweiten Spieler wird ein \enquote{Remote-Puppet}
    (ein lokal dargestellter, aber vom Netzwerk gesteuerter Würfel) geführt.}
    \eintrag{LevelManager}{Die Level-Generierung wurde serialisier- und deserialisierbar gemacht
    (\texttt{serialize\_level()}/\texttt{apply\_remote\_level()}), damit der Host das exakte
    Level-Layout überträgt, statt es beim Secondary separat und ggf. abweichend neu zu würfeln.}
    \eintrag{Player / Dash}{Die Kollisionslogik wurde erweitert, damit der Würfel des jeweils
    anderen Spielers wie ein Hindernis behandelt wird (siehe Kapitel~\ref{cha:coop-konzept}).}
    \eintrag{UI / MainMenu}{Ein neues Koop-Panel kam hinzu (Hosten/Beitreten, Namenswahl,
    Geräteliste per Discovery).}
    \eintrag{Enemy-Sync}{Der Gegnerzustand wird über eine stabile \texttt{sync\_id} statt über den
    Array-Index abgeglichen, damit z.B. eine vom Bomber mitten im Zug neu erzeugte Bombe auf beiden
    Geräten korrekt mitsynchronisiert wird.}
\end{infoblock}

% \section{Tests für den Koop-Modus}

% Zwei Regressionstests sichern die Netzwerk-Grundlagen ab: \texttt{\_test\_network\_manager\_cleanup}
% prüft, dass Sockets und Ports nach \texttt{stop()} sauber freigegeben werden und nicht mit einer
% folgenden Verbindung kollidieren; \texttt{\_test\_level\_serialization\_roundtrip} prüft, dass ein
% Level nach Serialisierung und Deserialisierung ohne Informationsverlust identisch rekonstruiert wird.


\chapter{Reflexion}
\label{cha:reflexion}

Besonders gut funktioniert hat die Kern-Mechanik des Würfels: Die in Kapitel~\ref{cha:implementierung}
beschriebene Trennung von visueller Rotation (Pivot-Verschiebung, Tween) und logischer
Raster-Positionsbestimmung ließ sich sauber umsetzen und musste im weiteren Projektverlauf nicht
grundlegend überarbeitet werden. Auch die modulare Architektur aus Kapitel~\ref{cha:architektur} --
Würfelseiten und Gegner als austauschbare Szenen mit fester Schnittstelle -- hat sich bewährt: Neue
Fähigkeiten und Gegnertypen ließen sich ergänzen, ohne bestehende Systeme anfassen zu müssen, was sich
später auch bei der Koop-Erweiterung als hilfreich erwies.

Die größte Herausforderung war, dass Godot für mich zu Projektbeginn eine vollständig unbekannte Engine
war (siehe KI-Nutzungserklärung) -- Node-/Szenen-System, Autoloads und Signals mussten parallel zur
eigentlichen Konzeption erlernt werden. Hinzu kam der Zeitdruck durch die in
Kapitel~\ref{cha:coop-konzept} beschriebene, kurzfristige Entscheidung für den Mehrspieler-Modus: Die
Umsetzung in der verbleibenden Zeit vor der Abschlusspräsentation ließ wenig Spielraum für Polish.
Auch das Balancing der Gegner -- ein Schwierigkeitsgefühl, das fordert, aber fair bleibt -- erwies sich
als schwieriger als gedacht und wurde iterativ über mehrere Testrunden angepasst.

Rückblickend würde ich den Mehrspieler-Modus bereits bei der ursprünglichen Konzeption zumindest
mitdenken, auch wenn er am Ende nicht von Anfang an umgesetzt worden wäre -- das hätte die späteren
Anpassungen an GameManager und LevelManager (siehe Kapitel~\ref{cha:coop-architektur}) vermutlich
etwas erleichtert. Am grundsätzlichen Vorgehen -- erst ein stabiles Singleplayer-Fundament, dann
gezielte Erweiterung -- würde ich aber nichts Wesentliches ändern.

Der größte Lernerfolg liegt entsprechend in zwei Bereichen: zum einen im Godot-spezifischen
Grundlagenwissen als Einsteiger, zum anderen im Umgang mit Architektur-Entscheidungen bei wachsendem
Projektumfang -- also darin, bestehende Systeme gezielt zu erweitern, statt sie bei neuen Anforderungen
neu zu bauen.


\appendix

%---
\printbibliography[heading=bibintoc]

\end{document}